% sections/external_experiments.tex - Section 4.4: the rule-based experiments as payment-path treatments. Ledger:
% external/experiments_paths.csv, external/experiments_estimates.csv; readings: Code/local/experiments_ledger.R,
% Code/local/experiments_rho_subgroups.R. Every number under the window aggregation of eq:window at the marginal-utility profile, with
% survival from the study's own exit hazards (revised 2026-09-16; the constant lambda = 0.2 is the sensitivity).
\subsection{Application to Rule-Based Experiments}\label{sec:quasi}

Proposition \ref{prop:curveT} states what an experiment must do to identify the discount curve: change the payment owed at a known date, by a known amount, for reasons unrelated to the borrower. I have identified seven prior studies that do exactly this, three of them with a pair of arms in one population. Each treatment is recorded as a series $\{(X_s, s)\}$, the change in dollars owed at month $s$ relative to the study's control (Figure \ref{fig:ledger_paths}, Appendix \ref{appendix:experiments}), and classified by the sign pattern of the series. A \emph{relief} has $X_s \le 0$ at every date; a \emph{rescheduling} has $X_s < 0$ at some dates and $X_s > 0$ at others. The distinction matters because the value of a rescheduling to a borrower depends on the weights the formula assigns to the dates it moves payments between: the discount function, survival, and the premium on the payment due now. The model's first-order default response to a treatment is $\phi \sum_s w_s X_s$ with $w_1 = 1$ and $w_s = D(s-1)\,Q_s\,\lambda_s$ for $s \ge 2$ (equation \eqref{eq:weights}), aggregated over the decision months of the study's outcome window as in equation \eqref{eq:window}, so within one population, where $\phi$ is common, the ratio of two arms' effects equals the ratio of their weighted sums and identifies $\rho$ given survival and the marginal-utility profile. Table \ref{tab:experiments} in Appendix \ref{appendix:experiments} lists every treatment and the test of each pair.

Two inputs are fixed before any pair is examined. Survival is the study's own: $Q_s = (1-h)^{s-1}$ with $h$ the monthly hazard of leaving the contract, the sum of the study's re-default or filing hazard, scaled by the probability that a spell of distress is not cured (0.795, from Appendix \ref{appendix:measuredrho}), and the study's payoff hazard. The marginal-utility profile is equation \eqref{eq:lambda_p} at the study population's own baseline rate and payment share, with its constant-income limit at the $0.2$ that Indarte's pair identifies below; a constant $0.2$ at every later date is the sensitivity. Every outcome is a window outcome, bankruptcy or delinquency within the study's horizon, and the test aggregates decisions over the window with the study's own event timing, as equation \eqref{eq:window} does for the bureau; the date-1 formula is reported beside it. Each pair is tested on a grid of the annual rate from 0 to 3 in steps of 0.01 by the joint $\chi^2(1)$ distance of the two reported effects from the model's line. Three reference rates are tested throughout: the 2 percent social discount rate of the 2023 revision of Circular A-4 \citep{the_white_house_biden-harris_2023}; the paper's estimate of $0.11$ (Section \ref{sec:rate_people}); and the laboratory literature's median of 0.22 per year \citep{cohen_measuring_2020}.

\citet{dobbie_targeted_2020} randomize the terms of debt-management plans, and this pair is the paper's check on the sign and the horizon of the far-payment response. The write-down arm forgives the last ten payments of a 52-month plan, a relief worth \$4,302 that arrives 3.5 to 4.3 years out; the payment-reduction arm lowers the payment by \$92 a month for 52 months and then adds 18 months of payments with \$1,645 more in fees, a rescheduling whose undiscounted value is negative. The write-down lowers five-year bankruptcy by $3.0$ percentage points (standard error $0.9$); the rescheduling raises it by $2.3$ ($1.1$). The model predicts both signs when the rescheduling's weighted value is negative, which, aggregated over the five-year window, holds at every rate below $1.4$: a filing in the third year of the window was decided with the forgiven payments one to two years ahead rather than four, and averaging over the window's decision months moves both arms' weighted values slowly with the rate. The joint test fits best at $\rho = 0$, and the rates consistent with the pair are $0$ to $1.20$ at the five percent level and $0$ to $2.84$ at the one percent level; the $p$-values at the three reference rates are $0.93$ at $0.02$, $0.73$ at $0.11$ and $0.54$ at $0.22$. What the pair rejects is the upper end of the grid, where the model predicts no response to the write-down and the observed response, three standard errors from zero, would be chance ($p = 0.009$ at $\rho = 3$), and with it a zero or negative weight on a payment four years out: that is the sign check, and the estimate of Section \ref{sec:rate_people} passes it. Under a constant premium of $0.2$ the sets end at $0.99$ and $1.97$; with the two arms' sampling errors correlated at $0.5$, their covariance being unreported, at $1.24$ and $2.24$; with survival set to one at $1.21$ and $2.85$, so the cure rate matters little. Read at a single decision at enrollment, the date-1 formula that a two-period model suggests, the same pair accepts only $0$ to $0.03$ at five percent and $0$ to $0.08$ at one percent; that formula assigns every filing in the window to a decision taken when the forgiven payments were four years away, and it overstates what the experiment says about the level. Figure \ref{fig:ds_pvalue} in Appendix \ref{appendix:figures} shows the whole curve. The timing corroborates the sign: the write-down's effect is concentrated in the first three years, before any forgiven payment falls due, as a forward-looking response to a far relief should be.

\citet{ganong_liquidity_2020} study HAMP modifications with two regression discontinuities, and this pair is the check at a long horizon. At the Treasury model's cutoff, borrowers receive \$31,000 more principal forgiveness with payments held at the same 31 percent of income for five years, so the relief arrives entirely after month 60; the effect on ninety-day delinquency over three years is $+1.2$ percentage points with a standard error of $3.2$. At the 31-percent payment-to-income cutoff, private modifications cut payments 19 percentage points more than HAMP for 22 years and then charge more for 18 years, a rescheduling with a small net relief; the effect on delinquency within two years is $-0.38$ points ($0.08$) per percent of payment reduction, $-7.2$ points ($1.5$) at the 19-point first stage. The forgiveness arm's standard error is three times its effect, so with the study's exit hazards no rate on the grid is rejected at either level, and the $p$-values at the reference rates are $0.25$, $0.40$ and $0.50$. Survival is what makes the pair interpretable at all: with $Q_s = 1$ the added payments after year 22 outweigh the deeper cut, so the rescheduling arm is a net obligation at rates below $0.06$; with the study's hazards, under which more than half of the borrowers have left the modified loan by the sixth year, the arm is a net relief at every rate. The pair contributes a 26-year horizon and no precision.

\citet{indarte_moral_2023} identifies two responses in household bankruptcy: filing falls by $0.21$ points ($0.06$) per \$1,000 of annual payment reduction at an adjustable-rate reset that lasts twelve months, and by $0.019$ points ($0.002$) per \$1,000 of seizable home equity at the filing date. The first is a relief within the year. The second is not a payment at any date; it is a change in the value of default realized at the decision, and its weight is the marginal utility of protected, illiquid wealth to a household at the filing margin. The pair therefore compares currencies rather than dates: its ratio is the marginal utility of an illiquid dollar relative to a liquid one for a household at the bankruptcy margin, and under the assumption of Corollary \ref{cor:separate}, that this is the marginal utility of a repayer's later dollar relative to the marginal defaulter's, it measures $\lambda$ rather than $\rho$. With the twelve months of relief weighted at the liquid dollar's weight, the observed ratio crosses the model's line at $\lambda \approx 0.2$ at every rate on the grid (Figure \ref{fig:ledger_rho}, bottom panel, Appendix \ref{appendix:experiments}), and that value is the one used throughout.

\citet{fiorin_moratorium_2026} offer delinquent consumer borrowers in India a three-month payment deferral with a processing fee, the deferred payments falling due after the original maturity: a rescheduling whose undiscounted value is the negative of the fee. When the offer is framed as regulatory guidance, it moves full repayment by $+0.3$ points ($1.3$) relative to a reminder call, indistinguishable from zero, which is what a rescheduling should do at a low rate. When the same offer is framed as the lender's own initiative it raises repayment by $3.3$ points ($1.4$); the difference between the two framings is not a difference in payment paths, so the model of this paper has nothing to say about it.

The experiments therefore fix the sign of the far-payment response and its horizon and leave its level to the bureau: over a four-year horizon the best fit is zero and every rate below $1.2$ is consistent with the pair at the five percent level; over a 26-year horizon nothing is measured. Section \ref{sec:mapping} states why: a response to a payment four years out identifies the rate less precisely than a response at the horizon of an auto contract, and an outcome window of years spreads the decision dates over which the far payment's weight is averaged. Turned into posterior weights under a flat prior, the results depend on the prior's support in a way that the $p$-value curve does not, and the mean of a posterior is not the right summary of this evidence; Appendix \ref{appendix:experiments} reports the posterior summaries by pair and support.

\citet{dobbie_targeted_2020} estimate both arms in one regression, so every subgroup they report is a pair, and the same test can be run within subgroups: by baseline debt-to-income above and below the median, by baseline credit score above and below the median, and by employment at baseline (their Table 8 and Appendix Table A9). Table \ref{tab:rho_subgroups} in Appendix \ref{appendix:experiments} reports the results as $p$-value curves rather than posterior sets, because the subgroup posteriors are bimodal, with a peak at zero and a second region of weight at high rates where the model predicts no response to the write-down and the subgroup's less precise write-down estimate is closer to zero in standard-error units. In every subgroup the best-fitting rate lies between zero and $0.20$, no reference rate is rejected at the five percent level, and the rates consistent at that level extend from zero to between $1.2$ (the pooled pair) and the top of the grid. The ratios of the two arms' effects do not differ between the halves of any split ($p = 0.64$ for debt-to-income, $0.84$ for credit score, $0.21$ for employment), and a common rate for the two halves is not rejected against separate ones ($p = 0.80$, $0.89$ and $0.64$). Within the precision of the experiment the rate over four years is the same for prime and subprime borrowers and for high- and low-leverage borrowers: a direct test of heterogeneity in the discount rate across the axes along which default rates differ most, and it finds none. Figure \ref{fig:rho_subgroups} there shows the curves.



The far payments in these designs are weighted by the probability of still facing them, and the two pairs show the two ways that weight enters. In the Dobbie--Song plan the exit hazard is low and the horizon short, so survival moves the sets little. In the Ganong--Noel modification the exit hazard is high and the horizon long, and survival decides the sign of the rescheduling arm's value at low rates. Section \ref{sec:autodata} applies the same weights to the credit-bureau data, where the payment margin can be estimated at population scale and the term margin is a different object.
